\documentclass[11pt,a4paper]{article}
\usepackage[utf8]{inputenc}
\usepackage[spanish]{babel}
\usepackage{amsmath, amssymb}
\usepackage{geometry}
\geometry{top=2.5cm, bottom=2.5cm, left=2.5cm, right=2.5cm}

\begin{document}

\section*{Serie de ejercicios seleccionados para preparar examen de los capítulos 1, 2 y 3}

\begin{itemize}
    \item[\textbf{1.}] Un protón, que es el núcleo de un átomo de hidrógeno, se representa como una esfera con un diámetro de $2.4 \text{ fm}$ y una masa de $1.67 \times 10^{-27} \text{ kg}$. (a) Determine la densidad del protón. (b) Establezca cómo se compara su respuesta del inciso (a) con la densidad del osmio.
    \vspace*{6cm}

    \item[\textbf{2.}] ¿Qué masa se requiere de un material con densidad $\rho$ para hacer un cascarón esférico hueco que tiene radio interior $r_1$ y radio exterior $r_2$?
    \vspace*{5cm}

    \item[\textbf{3.}] Suponga que la ecuación $x = At^3 + Bt$ describe el movimiento de un objeto en particular, siendo $x$ la dimensión de la longitud y $t$ la dimensión de tiempo. (a) Determine las dimensiones de las constantes $A$ y $B$. (b) Determine las dimensiones de la derivada $dx/dt = 3At^2 + B$.
    \vspace*{6cm}

    \item[\textbf{4.}] Sea $\rho_{Al}$ la densidad del aluminio y $\rho_{Fe}$ la del hierro. Encuentre el radio de una esfera de aluminio sólida que equilibra una esfera de hierro sólida de radio $r_{Fe}$ sobre una balanza de brazos iguales.
    \vspace*{5cm}
\end{itemize}

\newpage

\begin{itemize}
    \item[\textbf{5.}] Un auditorio mide $40.0 \text{ m} \times 20.0 \text{ m} \times 12.0 \text{ m}$. La densidad del aire es $1.20 \text{ kg/m}^3$. ¿Cuáles son (a) el volumen de la habitación en pies cúbicos y (b) el peso en libras del aire en la habitación?
    \vspace*{6cm}

    \item[\textbf{6.}] La densidad promedio del planeta Urano es $1.27 \times 10^3 \text{ kg/m}^3$. La proporción de la masa de Neptuno con la de Urano es $1.19$. La proporción del radio de Neptuno con el de Urano es $0.969$. Encuentre la densidad promedio de Neptuno.
    \vspace*{6cm}

    \item[\textbf{7.}] (a) ¿Cuál es el orden de magnitud del número de microorganismos en el tracto intestinal humano? Una escala de longitud bacteriana típica es de $10^{-6} \text{ m}$. Estime el volumen intestinal y suponga que $1\%$ del mismo está ocupada por una bacteria. (b) ¿El número de bacterias indica si estas son beneficiosas, peligrosas o neutras para el cuerpo humano? ¿Qué funciones podrían tener?
    \vspace*{6cm}

    \item[\textbf{8.}] Una mujer se encuentra a una distancia horizontal $x$ de una montaña y mide el ángulo $\theta$ de elevación de la cima sobre la horizontal. Después de caminar acercándose a la montaña una distancia $d$ a nivel del suelo, encuentra que el ángulo es $\phi$. Encuentre una ecuación general para la altura $y$ de la montaña en términos de $d$, $\phi$ y $\theta$; ignore la altura de los ojos al suelo.
    \vspace*{7cm}
\end{itemize}

\newpage

\begin{itemize}
    \item[\textbf{9.}] Una partícula se mueve de acuerdo con la ecuación $x = 10t^2$, donde $x$ está en metros y $t$ en segundos. (a) Encuentre la velocidad promedio para el intervalo de tiempo de $2.00 \text{ s}$ a $3.00 \text{ s}$. (b) Encuentre la velocidad promedio para el intervalo de tiempo de $2.00 \text{ s}$ a $2.10 \text{ s}$.
    \vspace*{6cm}

    \item[\textbf{10.}] Un electrón en un tubo de rayos catódicos acelera uniformemente desde una rapidez de $2.00 \times 10^4 \text{ m/s}$ a $6.00 \times 10^6 \text{ m/s}$ en $1.50 \text{ cm}$. (a) ¿En qué intervalo de tiempo el electrón recorre estos $1.50 \text{ cm}$? (b) ¿Cuál es su aceleración?
    \vspace*{6cm}

    \item[\textbf{11.}] Un atacante en la base de la pared de un castillo de $3.65 \text{ m}$ de alto lanza una roca recta hacia arriba con una rapidez de $7.40 \text{ m/s}$ desde una altura de $1.55 \text{ m}$ sobre el suelo. (a) ¿La roca llegará a lo alto de la pared? (b) Si es así, ¿cuál es su rapidez en lo alto? Si no, ¿qué rapidez inicial debe tener para llegar a lo alto? (c) Encuentre el cambio en rapidez de una roca lanzada recta hacia abajo desde lo alto de la pared con una rapidez inicial de $7.40 \text{ m/s}$ y que se mueve entre los mismos dos puntos. (d) ¿El cambio en la rapidez de la roca que se mueve hacia abajo concuerda con la magnitud del cambio de rapidez de la roca que se mueve hacia arriba entre las mismas elevaciones? Explique físicamente por qué sí o por qué no concuerda.
    \vspace*{8cm}

    \item[\textbf{12.}] Se lanza una pelota hacia arriba desde el suelo con una rapidez inicial de $25 \text{ m/s}$; en el mismo instante, se deja caer otra pelota desde un edificio de $15 \text{ m}$ de altura. ¿Después de cuánto tiempo estarán las pelotas a la misma altura por encima del suelo?
    \vspace*{6cm}
\end{itemize}

\newpage

\begin{itemize}
    \item[\textbf{13.}] Al tiempo $t = 0$, un estudiante lanza un juego de llaves verticalmente hacia arriba a su hermana de fraternidad, que está en una ventana a una distancia $h$ arriba. La segunda estudiante atrapa las llaves al tiempo $t$. (a) ¿Con qué velocidad inicial se lanzaron las llaves? (b) ¿Cuál era la velocidad de las llaves justo antes de que fueran capturadas?
    \vspace*{6cm}

    \item[\textbf{14.}] Una mujer informó haber caído $144$ metros del piso 17 de un edificio, quedando sobre una caja de metal del ventilador que aplastó a una profundidad de $18.0 \text{ pulg}$. Ella sufrió solo heridas leves. Ignorando la resistencia del aire, calcule (a) la rapidez de la mujer justo antes de que chocó con el ventilador y (b) su aceleración promedio mientras está en contacto con la caja. (c) Modele su aceleración como constante, calcule el intervalo de tiempo que se tardó en aplastar la caja.
    \vspace*{6cm}

    \item[\textbf{15.}] Usted está sentado en su auto en reposo en un semáforo con un ciclista en reposo junto a usted en el carril contiguo de bicicletas. Tan pronto como el semáforo se pone verde, un automóvil acelera a partir del reposo a $50.0 \text{ mi/h}$ con una aceleración constante de $9.00 \text{ mi/h/s}$ y después se mueve con una rapidez constante de $50.0 \text{ mi/h}$. Al mismo tiempo el ciclista acelera desde el reposo hasta $20.0 \text{ mi/h}$ con una aceleración constante de $13.0 \text{ mi/h/s}$ y después se mueve con rapidez constante de $20.00 \text{ mi/h}$. (a) ¿Durante qué intervalo de tiempo está la bicicleta por delante del automóvil? (b) ¿Para qué distancia máxima la bicicleta está delante del automóvil?
    \vspace*{8cm}

    \item[\textbf{16.}] El vector $\vec{A}$ tiene una magnitud de $29$ unidades y apunta en la dirección $y$ positiva. Cuando el vector $\vec{B}$ se suma al vector $\vec{A}$, el vector resultante $\vec{A} + \vec{B}$ apunta en dirección $y$ negativa con una magnitud de $14$ unidades. Encuentre la magnitud y la dirección de $\vec{B}$.
    \vspace*{6cm}
\end{itemize}

\newpage

\begin{itemize}
    \item[\textbf{17.}] Un carro de montaña rusa se mueve $200 \text{ pies}$ horizontalmente y luego se eleva $135 \text{ pies}$ en un ángulo de $30.0^\circ$ sobre la horizontal. A continuación, viaja $135 \text{ pies}$ en un ángulo de $40.0^\circ$ hacia abajo. ¿Cuál es su desplazamiento desde su punto de partida? Use técnicas gráficas.
    \vspace*{8cm}

    \item[\textbf{18.}] Una persona camina $25.0^\circ$ al noreste durante $3.10 \text{ km}$. ¿Qué distancia tendría que caminar hacia el norte y hacia el este para llegar a la misma posición?
    \vspace*{6cm}

    \item[\textbf{19.}] El vector $\vec{A}$ tiene componentes $x$ y $y$ de $-8.70 \text{ cm}$ y $15.0 \text{ cm}$, respectivamente; el vector $\vec{B}$ tiene componentes $x$ y $y$ de $13.2 \text{ cm}$ y $-6.60 \text{ cm}$, respectivamente. Si $\vec{A} - \vec{B} + 3\vec{C} = 0$ ¿cuáles son los componentes de $\vec{C}$?
    \vspace*{6cm}

    \item[\textbf{20.}] En la figura a continuación se muestran tres vectores desplazamiento de una pelota de croquet, donde $|\vec{A}| = 20.0$ unidades, $|\vec{B}| = 40.0$ unidades y $|\vec{C}| = 30.0$ unidades. Encuentre (a) la resultante en notación de vectores unitarios y (b) la magnitud y dirección del desplazamiento resultante.
    \vspace*{6cm}
\end{itemize}

\newpage

\begin{itemize}
    \item[\textbf{21.}] Una chica que entrega periódicos cubre su ruta viajando $3.00$ cuadras al oeste, $4.00$ cuadras al norte, y luego $6.00$ cuadras al este. (a) ¿Cuál es su desplazamiento resultante? (b) ¿Cuál es la distancia total que viaja?
    \vspace*{6cm}

    \item[\textbf{22.}] Un transbordador lleva turistas entre tres islas. Navega de la primera isla a la segunda, a $4.76 \text{ km}$ de distancia, en una dirección $37.0^\circ$ al noreste. Luego navega de la segunda isla a la tercera en una dirección de $69.0^\circ$ al noroeste. Por último, regresa a la primera isla y navega en una dirección $28.0^\circ$ al sureste. Calcule la distancia entre (a) la segunda y tercera islas y (b) la primera y tercera islas.
    \vspace*{6cm}

    \item[\textbf{23.}] Los vectores $\vec{A}$ y $\vec{B}$ tienen magnitudes iguales de $5.00$. La suma de $\vec{A}$ y $\vec{B}$ es el vector $6.00\hat{j}$. Determine el ángulo entre $\vec{A}$ y $\vec{B}$.
    \vspace*{6cm}
\end{itemize}

\end{document}
